\section{Introduction}
\label{sec:introduction}

Blended Wing Body (BWB) and tailless flying-wing configurations provide distinct aerodynamic and structural advantages for long-range autonomous Unmanned Aerial Vehicle (UAV) applications~\cite{liebeck2004design,qin2004aerodynamic,okonkwo2016review}. By integrating the fuselage and lifting surfaces into a continuous lifting body, these configurations reduce the parasite drag, wetted area and wing-body interference associated with conventional cylindrical fuselage and discrete empennage architectures~\cite{kroo2001innovative,torenbeek2013advanced}. The deep center-body section also yields substantial internal volume, providing space for high-density battery packaging and operational payloads. Multidisciplinary design optimization (MDO) frameworks are consequently central to resolving the coupled aerodynamic, structural, and weight tradeoffs that characterize these vehicles~\cite{martins2013multidisciplinary,kenway2014multipoint,brooks2018benchmark}.

Tailless BWB configurations must trim and remain statically stable without a tail, through sweep, washout and reflexed sections~\cite{nickel1994tailless,nelson1998flight}. Their static margin depends on where the battery sits relative to a neutral point that the planform sets, and low-order potential-flow models evaluate any planform that a parameterization proposes, whether or not it is physically sensible~\cite{katz2001low}. The feasible region is therefore small and its boundary depends on the fidelity of the model, which is the setting of the present study; whether a geometric safeguard on the parameterization improves a search is tested in Section~\ref{sec:loft}.

Surrogate-assisted optimization methods, including Efficient Global Optimization (EGO)~\cite{jones1998efficient}, Kriging and other surrogate modeling frameworks~\cite{forrester2008engineering,koziel2013variable,bouhlel2019python}, and active learning algorithms~\cite{settles2009active,shahriari2016human}, are standard tools for reducing computational expense in aircraft conceptual design. Recent investigations in BWB optimization have applied adaptive Kriging~\cite{suo2026kriging}, hybrid geometric parameterizations~\cite{wang2026hybrid}, aerodynamic shape optimization~\cite{lyu2014aerodynamic,reist2013aerodynamic}, aerostructural coupling~\cite{jasa2018openaerostruct}, and variable-fidelity models~\cite{koziel2013variable,koziel2013surrogate}.

Several lines of prior work are closest to the present study, and each differs from it in scope. Hakim et al.\ build a Gaussian-process static-stability design framework for tailless BWB mini-UAVs from XFLR5 samples and Latin hypercube sampling~\cite{hakim2026integrated}; the present study instead compares search strategies under an equal evaluation budget and audits the stability constraint against Euler CFD. Mader and Martins treat stability-constrained shape optimization of flying wings with gradient-based methods~\cite{mader2013stability}, whereas the present problem is gradient-free and parameterized by a 14-variable conceptual vector. Wang et al.\ optimize a tailless UAV for endurance under center-of-gravity location uncertainty~\cite{wang2015tailless}; the battery-placement sweep in Section~\ref{sec:dynamic_balance} examines the same sensitivity for the best design of the original campaign. Finally, Plaban et al.\ report that derivative-free optimization is competitive for aerodynamic design in moderate dimensions~\cite{plaban2025derivative}, which is consistent with the equal-budget results reported here, and Cordelier et al.\ discuss multi-fidelity constrained Bayesian optimization for aircraft design and the limits of probability-of-feasibility weighting~\cite{cordelier2026multifidelity}, which is the weighting used by the active-learning loop studied here. Two recent studies of small BWB UAVs use surrogate-assisted search directly. Naseem et al.\ train a Gaussian-process surrogate on CFD results for a one-meter BWB UAV, optimize it with Bayesian optimization, and build and wind-tunnel test a prototype~\cite{naseem2026blended}; their study validates the final design experimentally, which the present one does not, and the present study asks a question theirs does not, namely how much of the improvement a surrogate-free search would have found with the same number of evaluations. Suo et al.\ optimize the range of a BWB UAV in eight variables under a budget of 80 simulations and compare a one-shot Kriging model with two adaptive Kriging strategies over 15 randomized trials each~\cite{suo2026kriging}; the median validated ranges of the three strategies lie within 0.2\,km of each other, and a sensitivity analysis finds the fuel mass to be the most influential variable. That study replicates a comparison between surrogate strategies; the present one adds controls without any surrogate, a stability constraint, and an ablation of how feasibility is handled.

Nevertheless, three critical methodological questions remain under-explored in the surrogate-assisted aircraft design literature:
\begin{itemize}
    \item \textit{Equal-Budget Attribution:} When a surrogate-guided active learning loop discovers an improved configuration, how much of that gain is attributable to the surrogate and how much to the localized search budget or the choice of seed configurations? Rigorous empirical protocols in machine learning emphasize equal-budget baselines~\cite{bergstra2012random} and multi-seed replications~\cite{henderson2018deep}, yet surrogate-assisted aerospace MDO studies frequently compare active learning outcomes solely against initial space-filling pools rather than direct local search baselines operating under identical evaluation counts.
    \item \textit{Bounds and Constraint Fidelity versus Search Strategy:} How much of the achievable performance is set by the variable bounds and by the fidelity of the stability constraint, compared with the choice of search algorithm, and what does a geometric safeguard such as a chord clamp contribute?
    \item \textit{Low-Order Aerodynamic Fidelity and Consistency:} To what degree can a low-order potential-flow pipeline incorporating empirical strip drag be trusted, both in design ranking and absolute performance, when it operates as the sole evaluation engine for the optimizer?
\end{itemize}

This paper addresses these questions for a 1.5\,m-span, 8\,kg electric tailless BWB UAV parameterized by 14 continuous design variables. Its contributions are the following.
\begin{enumerate}
    \item \textbf{An equal-budget benchmark of search strategies with a surrogate-free control matched to the surrogate loop.} On a shared 9,800-design pool, eight strategies receive 287 further evaluations each over 10 to 20 seeds: an MLP-ensemble active-learning loop with a feasibility-classifier weight, the same loop with a washout repair, TuRBO~\cite{eriksson2019turbo}, an adaptive local search that uses the same batch structure and parent set as the loop but random proposals, two static local searches, and Latin hypercube and uniform sampling (Sections~\ref{sec:protocol} and~\ref{sec:budget_results}). Comparisons use Welch and Mann--Whitney tests, bootstrap intervals and Holm correction, and the leading strategies are repeated from pools of 288 to 3{,}000 designs (Section~\ref{sec:pool_size}).
    \item \textbf{A factorial analysis of feasibility handling.} The classifier weight (sigmoid, hard gate, none) is crossed with the washout repair, and the weight is also applied to TuRBO, which separates what is necessary for the loop from what is incidental (Section~\ref{sec:feas_ablation}).
    \item \textbf{A stability constraint calibrated against Euler CFD.} SU2 Euler solutions for \NcalN\ designs give a linear map from the lattice static margin to the Euler margin; the map is then checked on \NoosN\ held-out designs and every run is re-ranked under the calibrated constraint (Section~\ref{sec:sm_calibration}).
    \item \textbf{Bound prices and loft audit.} Single-bound relaxations measure the range cost of each active bound (Section~\ref{sec:bound_prices}), and a controlled comparison shows what the chord clamp of the planform loft does and does not do (Section~\ref{sec:geometric_manifold}).
    \item \textbf{A code-to-text audit with disclosed corrections.} Every equation, constant and algorithmic statement in the paper was checked against the code that produced the results. The audit found an energy-basis inconsistency in two baselines, an asymmetry in constraint handling between strategies, and an incorrect monotonicity claim for the loft. All three are corrected, the earlier runs are retained for comparison (Table~\ref{tab:asrun}), and the conclusions are stated on the corrected results.
\end{enumerate}
The evaluator is checked by panel refinement, a second aerodynamic model (AeroBuildup), learned section polars (NeuralFoil, which AeroBuildup also uses), and Euler CFD (Section~\ref{sec:validation}). Table~\ref{tab:claims} states the evidence and the limit of each claim.
\begin{table}[htbp]
\centering\small
\caption{Claims, evidence and scope limits.}
\label{tab:claims}
\resizebox{\linewidth}{!}{%
\begin{tabular}{@{}p{0.22\linewidth}p{0.42\linewidth}p{0.36\linewidth}@{}}
\toprule
\textbf{Claim} & \textbf{Evidence} & \textbf{Scope limit} \\
\midrule
TuRBO is the strongest strategy; the active-learning loop needs a washout repair to exceed a surrogate-free adaptive local search & \NvNTurbo\ TuRBO, \NvNALR\ AL-R, \NvNAL\ AL and \NvNAdapt\ adaptive-local runs on one pool with 287 evaluations each (Tables~\ref{tab:budget} and~\ref{tab:tests}); ordering repeated at four pool sizes and five budgets (Tables~\ref{tab:pool} and~\ref{tab:anytime}) & One 14-variable problem and one evaluator; independent pools were tested at one size only. The lead of TuRBO over AL-R is \NvPctTurbovALR\% (Holm $p \NvHTurbovALR$) and shrinks, with all gaps, as the budget grows. \\
\addlinespace
The classifier weight is necessary for the loop and unhelpful for TuRBO & Factorial runs, ten seeds per cell (Table~\ref{tab:factorial}) & Applies to a global acquisition on surrogates fitted to feasible designs only. \\
\addlinespace
The lattice static margin is low by a nearly constant offset & Euler CFD on \NcalN\ designs, $R^2 = \NcalRsq$ (Fig.~\ref{fig:sm_calibration}) & Euler is inviscid; the true margin needs RANS, wind-tunnel or flight data. \\
\addlinespace
The battery-mass bound is the only bound with a measurable price; the chord clamp is a validity guard & Single-bound relaxations, five seeds per bound (Table~\ref{tab:bound_prices}); clamp versus no clamp, ten paired seeds (Table~\ref{tab:loft}) & Five seeds leave gains of a few kilometers possible for the other bounds. The loft does not guarantee a monotone chord inside the center body. \\
\addlinespace
The evaluator ranks designs consistently with independent models & Panel refinement, AeroBuildup ($n=\NxN$), NeuralFoil, Euler CFD (Tables~\ref{tab:grid}, \ref{tab:visc} and~\ref{tab:cfd}) & No viscous three-dimensional solution and no experiment; absolute range is not validated. \\
\bottomrule
\end{tabular}%
}
\end{table}
